% slides/03_eeg_primer.tex
\begin{frame}{EEG 101}
  \begin{columns}[T]
    \begin{column}{0.50\textwidth}
      \textbf{Generación}
      \begin{itemize}
        \item Potenciales post-sinápticos en las dendritas apicales de neuronas piramidales
        \item \emph{¡También potenciales de acción!}
        \item Conducción por el volumen de LCR/cráneo/piel cabelluda
      \end{itemize}
      \vspace{0.8em}
      \textbf{Interpretación \emph{clásica} en bandas}
        \small
        \begin{tabular}{lll}
          \toprule
          Banda & Rango (Hz) & Correlato cognitivo \\
          \midrule
          $\delta$ & 0.5--4   & Sueño profundo, patología \\
          $\theta$ & 4--8     & Memoria, somnolencia \\
          $\alpha$ & 8--13    & Vigilia \\
          $\beta$  & 13--30   & Cognición activa \\
          $\gamma$ & 30--100+ & Fijación, percepción \\
          \bottomrule
        \end{tabular}
    \end{column}
    \begin{column}{0.5\textwidth}
    	\begin{block}{Referencia clave}
    		\small
    		Fernandez-Ruiz \textit{et al.} (2023) asocian las oscilaciones gamma con la operación de circuitos neuronales \cite{Fernandez-Ruiz2023}.
    	\end{block}
    \end{column}
  \end{columns}
\end{frame}
